\needspace{10\baselineskip}
\hypertarget{native-322-zusatz17}{}
\pdfbookmark[1]{17. Zodiakus auf der Himmelskugel}{zusatz-017}
\begin{center}
[17.]

Zodiakus auf der Himmelskugel.

28. Okt. 1847.
\end{center}

Für jeden Punkt der Himmelskugel gibt es entweder keine, eine oder mehrere Kombinationen, durch die der eine Planet vom anderen aus an diesem Punkt sichtbar ist, also eine bestimmte Anzahl von Lösungen. Beim Überschreiten der Grenzlinie ändert sich die Anzahl der Lösungen jedes Mal um 2. Solange man aber auf einem beliebigen Weg auf der Himmelskugel die Grenzlinie nicht trifft, bleibt die Zahl der Lösungen unverändert. Die Frage ist nun, wie man die Seite der Grenze mit zwei Lösungen mehr von der anderen Seite mit zwei Lösungen weniger unterscheidet. Dazu dienen die folgenden Überlegungen.

1. Es sei \(p\) ein Punkt der Planetenbahn, von der aus die Punkte der anderen Bahn auf die Himmelskugel projiziert werden, \(P\) der \SourceLeaf{323}entsprechende Punkt der anderen Bahn und \(\Pi\) der Punkt der Knotenlinie, an dem sich die Tangenten an die beiden Planetenbahnen in \(p\) und \(P\) schneiden. Ferner sei
\[
p\Pi=n,\qquad P\Pi=N,\qquad pP=\Delta,
\]
\(S\) ein anderer beliebiger fester Punkt der Knotenlinie (etwa die Sonne). Die drei Figuren sind in drei verschiedenen Ebenen vorzustellen: I in der Ebene durch \(\Pi pP\), II in der Ebene durch \(\Pi Sp\), III in der Ebene durch \(\Pi SP\). Ferner bezeichnen wir mit
\begin{quote}
\noindent \(p,P\) die Winkel der Ebenen II und III mit I,\\
\(q,Q\) die Winkel \(p\Pi S,P\Pi S\),\\
\(u,U\) die Winkel \(\Pi pP,\Pi Pp\),\\
\(r,R\) die Krümmungsradien der beiden Bahnen in \(p\) und \(P\).
\end{quote}

Nun seien \(p^{\prime},P^{\prime}\) zwei andere einander entsprechende Punkte und \(\Pi^{\prime}\) der zugehörige Punkt der Knotenlinie. Dabei wird \(\Pi\Pi^{\prime}=\delta\) als unendlich klein angenommen. Dann ergeben sich für die Lote von \(p^{\prime},P^{\prime}\) auf \(pP\):
\[
p^{*}p^{\prime}=\frac{\delta r\sin q\sin u}{n}\mathpunct{\raisebox{.45ex}{\normalfont,}}\qquad
P^{*}P^{\prime}=\frac{\delta R\sin Q\sin U}{N}\mathpunct{\cdot}
\]

Ferner seien \(p^{\prime\prime},P^{\prime\prime}\) zwei andere Punkte auf den Tangenten \(p\Pi,P\Pi\), sodass \(p^{\prime\prime}P^{\prime\prime}\) zu \(pP\) parallel und \(\frac{pp^{\prime\prime}}{p\Pi}=\frac{PP^{\prime\prime}}{P\Pi}=\omega\) unendlich klein ist. Eine durch \(p^{\prime\prime}P^{\prime\prime}\) senkrecht zu I gelegte Ebene treffe die Kurven in \(p^{\prime\prime\prime},P^{\prime\prime\prime}\), und \(p^{\mathrm{IV}},P^{\mathrm{IV}}\) seien die Fußpunkte der Lote von \(p^{\prime\prime\prime},P^{\prime\prime\prime}\) auf I. Dann gilt
\[
\begin{aligned}
p^{\prime\prime\prime}p^{\mathrm{IV}}&=p^{\prime\prime}p^{\prime\prime\prime}\mathord{.}\sin p=\frac{nn\omega\omega}{2r}\sin p\\[5pt]
P^{\prime\prime\prime}P^{\mathrm{IV}}&=P^{\prime\prime}P^{\prime\prime\prime}\mathord{.}\sin P=\frac{NN\omega\omega}{2R}\sin P\\[5pt]
p^{\prime\prime\prime}p^{\mathrm{IV}}\mathord{.}p^{*}p^{\prime}&=\frac{\delta\omega\omega n\sin p\sin q\sin u}{2}\\[5pt]
P^{\prime\prime\prime}P^{\mathrm{IV}}\mathord{.}P^{*}P^{\prime}&=\frac{\delta\omega\omega N\sin P\sin Q\sin U}{2}\mathpunct{\cdot}
\end{aligned}
\]

Man sieht aber leicht, dass \(n\sin u=N\sin U\) das Lot von \(\Pi\) auf \(pP\) ist. Ferner gilt \(\sin p\sin q=\sin P\sin Q={}\) Sinus des Winkels, den die Gerade \(\Pi S\) mit der Ebene I und folglich mit dem genannten Lot bildet. Daraus schließt man leicht, dass \(n\sin p\sin q\sin u=N\sin P\sin Q\sin U\) der kürzeste Abstand zwischen den Geraden \(pP,\Pi S\) ist.

\SourceLeaf{324}Das Endergebnis ist also
\[
p^{\prime\prime\prime}p^{\mathrm{IV}}\mathord{.}p^{*}p^{\prime}=P^{\prime\prime\prime}P^{\mathrm{IV}}\mathord{.}P^{*}P^{\prime}.
\]
Ist folglich \(p^{\prime\prime\prime}p^{\mathrm{IV}}>P^{\prime\prime\prime}P^{\mathrm{IV}}\), so gilt \(p^{*}p^{\prime}<P^{*}P^{\prime}\) und umgekehrt. Die erste Ungleichung zeigt an, dass eine zu \(p^{\prime\prime\prime}P^{\prime\prime\prime}\) parallele, von \(p\) aus gezogene Gerade und der Punkt \(S\) auf entgegengesetzten Seiten von I liegen. Die zweite zeigt an, dass die Projektionen der beiden Geraden \(p^{\prime}P^{\prime}\) und \(pP^{\prime}\) auf die Himmelskugel (wobei die zweite Gerade nur um eine unendlich kleine Größe 2\textsuperscript{ter} Ordnung von der Grenzlinie entfernt ist) auf derselben Seite der Projektion von \(pP\) liegen.

Daraus lässt sich das gesuchte Kriterium leicht ableiten und wie folgt formulieren:

[Sei] \(B\) [die] äußere Planetenbahn, von der Sonne aus auf die Himmelskugel bezogen. Dieser Großkreis teilt die Himmelskugel in zwei Hälften (1) und (2). Sei (1) die Hälfte, in der der heliozentrische Ort des inneren Planeten erscheint. Dann liegt die Gegend mit zwei Lösungen mehr in Bezug auf ein Stück \(LL^{\prime}\) der Grenze genauso wie (1) in Bezug auf \(B\), wenn auf \(LL^{\prime}\) und \(B\) einander entsprechende Richtungen als vorwärts gelten. Dabei wird vorausgesetzt, dass der Planet auf der äußeren Bahn beobachtet wird. Im entgegengesetzten Fall verhält es sich umgekehrt.

In Symbolen lässt sich das Kriterium noch allgemeiner ausdrücken:

Es seien \(l,l^{\prime}\) zwei einander unendlich nahe Punkte auf einer Grenze, \(p,p^{\prime}\) die entsprechenden heliozentrischen Örter des Planeten, von dem aus [beobachtet wird], \(P,P^{\prime}\) dagegen die des Planeten, der beobachtet wird, und \(a\) ein Punkt nahe \(ll^{\prime}\) auf der Seite mit zwei Lösungen mehr als auf der anderen. Wenn das Kongruenzzeichen \(\equiv\) die gleiche Orientierung [bezeichnet] und ein vorangestelltes \(-\)-Zeichen die entgegengesetzte Orientierung bedeutet, gilt:
\[
ll^{\prime}a\equiv PP^{\prime}p\equiv-pp^{\prime}P.
\]
\begin{center}\rule{.22\textwidth}{.3pt}\end{center}
